\begin{textsample}{POS Dim 1 – vem\_ed\_04 – Score 16.00 – vem\_ed\_04/t016.txt}  \label{ex:f1_pos_179}
BOAS PRÁTICAS Com o pé direito Fatec Ipiranga estreia nos \textbf{intercâmbios} virtuais com um projeto desenvolvido com a Jamestown Community College, ligada à State University of New York (JCC / SUNY).
As atividades \textbf{foram} conduzidas por José Carlos Barbosa Lopes, \textbf{que} leciona Inglês na Fatec Ipiranga, e Renee Funke, da área de Educação da JCC / SUNY.
No primeiro \textbf{encontro}, em 18 de setembro, os 26 alunos brasileiros perguntaram via Zoom sobre a experiência \textbf{profissional} da professora norte-americana, a adaptação às aulas online e a empregabilidade na pandemia.
Para preparar os estudantes, Lopes oferece oficinas de conversação em inglês – além das duas horas-aula semanais do \textbf{idioma} obrigatórias nos currículos das Fatecs, diferencial ante outras instituições de ensino \textbf{superior} tecnológico.
Segundo Lopes, “o projeto \textbf{é} importante para os alunos da Fatec Ipiranga, não só pelo aprimoramento da língua inglesa, mas também pelo contato com estudantes de outro país e a experiência de discutir assuntos de interesse global”.
Sobre o planejamento do \textbf{intercâmbio} virtual, o professor afirma \textbf{que} \textbf{foram} considerados os perfis do alunado das duas instituições, os objetivos das \textbf{disciplinas} e as temáticas afins para compor a colaboração.
Na fase “quebra-gelo”, os 40 participantes (no Brasil e nos EUA) elaboraram um vídeo pessoal, trocaram mensagens e definiram os assuntos mais \textbf{relevantes} para o grupo – e assim \textbf{foram} montadas as equipes mistas (Fatec Ipiranga e JCC / SUNY).
As ferramentas digitais usadas \textbf{foram} Zoom, Padlet, Facebook, WhatsApp e Teams.
Como a maioria dos 26 participantes da Fatec Ipiranga frequenta aulas e oficinas de Inglês, \textbf{foi} \textbf{possível} “discutir e avaliar nos grupos a \textbf{importância} de \textbf{fatores} como colaboração, relações interpessoais, aspectos culturais e recursos de negociação de sentido, \textbf{todos} fundamentais num projeto de \textbf{intercâmbio} internacional”, relata Lopes. “O fato de poderem participar de discussões em inglês sobre assuntos tão \textbf{relevantes} atualmente como pandemia, multiculturalidade, \textbf{questões} sociais e \textbf{políticas}, \textbf{foi} inspirador”.
Perguntado sobre \textbf{que} dica daria para professores \textbf{que} desejam \textbf{começar} um \textbf{intercâmbio} virtual, Lopes recomenda: “Creio \textbf{que} \textbf{seja} interessante \textbf{observar} as possibilidades de aperfeiçoamento \textbf{que} poderão \textbf{ser} agregadas ao curso, à formação dos alunos e à \textbf{disciplina}”.
Cabe lembrar \textbf{que} a equipe dos PCIs pode ajudar em \textbf{todos} esses passos.

% matched lemmas: começar, disciplina, encontro, fator, idioma, importância, intercâmbio, observar, político, possível, profissional, que, questão, relevante, ser, superior, todo
\end{textsample}
